\IfFileExists{stacks-project.cls}{%
\documentclass{stacks-project}
}{%
\documentclass{amsart}
}

% For dealing with references we use the comment environment
\usepackage{verbatim}
\newenvironment{reference}{\comment}{\endcomment}
%\newenvironment{reference}{}{}
\newenvironment{slogan}{\comment}{\endcomment}
\newenvironment{history}{\comment}{\endcomment}

% For commutative diagrams we use Xy-pic
\usepackage[all]{xy}

% We use 2cell for 2-commutative diagrams.
\xyoption{2cell}
\UseAllTwocells

% We use multicol for the list of chapters between chapters
\usepackage{multicol}

% This is generally recommended for better output
\usepackage{lmodern}
\usepackage[T1]{fontenc}

% For cross-file-references
\usepackage{xr-hyper}

% Package for hypertext links:
\usepackage{hyperref}

% For any local file, say "hello.tex" you want to link to please
% use \externaldocument[hello-]{hello}
\externaldocument[introduction-]{introduction}
\externaldocument[conventions-]{conventions}
\externaldocument[sets-]{sets}
\externaldocument[categories-]{categories}
\externaldocument[topology-]{topology}
\externaldocument[sheaves-]{sheaves}
\externaldocument[sites-]{sites}
\externaldocument[stacks-]{stacks}
\externaldocument[fields-]{fields}
\externaldocument[algebra-]{algebra}
\externaldocument[brauer-]{brauer}
\externaldocument[homology-]{homology}
\externaldocument[derived-]{derived}
\externaldocument[simplicial-]{simplicial}
\externaldocument[more-algebra-]{more-algebra}
\externaldocument[smoothing-]{smoothing}
\externaldocument[modules-]{modules}
\externaldocument[sites-modules-]{sites-modules}
\externaldocument[injectives-]{injectives}
\externaldocument[cohomology-]{cohomology}
\externaldocument[sites-cohomology-]{sites-cohomology}
\externaldocument[dga-]{dga}
\externaldocument[dpa-]{dpa}
\externaldocument[sdga-]{sdga}
\externaldocument[hypercovering-]{hypercovering}
\externaldocument[schemes-]{schemes}
\externaldocument[constructions-]{constructions}
\externaldocument[properties-]{properties}
\externaldocument[morphisms-]{morphisms}
\externaldocument[coherent-]{coherent}
\externaldocument[divisors-]{divisors}
\externaldocument[limits-]{limits}
\externaldocument[varieties-]{varieties}
\externaldocument[topologies-]{topologies}
\externaldocument[descent-]{descent}
\externaldocument[perfect-]{perfect}
\externaldocument[more-morphisms-]{more-morphisms}
\externaldocument[flat-]{flat}
\externaldocument[groupoids-]{groupoids}
\externaldocument[more-groupoids-]{more-groupoids}
\externaldocument[etale-]{etale}
\externaldocument[chow-]{chow}
\externaldocument[intersection-]{intersection}
\externaldocument[pic-]{pic}
\externaldocument[weil-]{weil}
\externaldocument[adequate-]{adequate}
\externaldocument[dualizing-]{dualizing}
\externaldocument[duality-]{duality}
\externaldocument[discriminant-]{discriminant}
\externaldocument[derham-]{derham}
\externaldocument[local-cohomology-]{local-cohomology}
\externaldocument[algebraization-]{algebraization}
\externaldocument[curves-]{curves}
\externaldocument[resolve-]{resolve}
\externaldocument[models-]{models}
\externaldocument[functors-]{functors}
\externaldocument[equiv-]{equiv}
\externaldocument[pione-]{pione}
\externaldocument[etale-cohomology-]{etale-cohomology}
\externaldocument[proetale-]{proetale}
\externaldocument[relative-cycles-]{relative-cycles}
\externaldocument[more-etale-]{more-etale}
\externaldocument[trace-]{trace}
\externaldocument[crystalline-]{crystalline}
\externaldocument[spaces-]{spaces}
\externaldocument[spaces-properties-]{spaces-properties}
\externaldocument[spaces-morphisms-]{spaces-morphisms}
\externaldocument[decent-spaces-]{decent-spaces}
\externaldocument[spaces-cohomology-]{spaces-cohomology}
\externaldocument[spaces-limits-]{spaces-limits}
\externaldocument[spaces-divisors-]{spaces-divisors}
\externaldocument[spaces-over-fields-]{spaces-over-fields}
\externaldocument[spaces-topologies-]{spaces-topologies}
\externaldocument[spaces-descent-]{spaces-descent}
\externaldocument[spaces-perfect-]{spaces-perfect}
\externaldocument[spaces-more-morphisms-]{spaces-more-morphisms}
\externaldocument[spaces-flat-]{spaces-flat}
\externaldocument[spaces-groupoids-]{spaces-groupoids}
\externaldocument[spaces-more-groupoids-]{spaces-more-groupoids}
\externaldocument[bootstrap-]{bootstrap}
\externaldocument[spaces-pushouts-]{spaces-pushouts}
\externaldocument[spaces-chow-]{spaces-chow}
\externaldocument[groupoids-quotients-]{groupoids-quotients}
\externaldocument[spaces-more-cohomology-]{spaces-more-cohomology}
\externaldocument[spaces-simplicial-]{spaces-simplicial}
\externaldocument[spaces-duality-]{spaces-duality}
\externaldocument[formal-spaces-]{formal-spaces}
\externaldocument[restricted-]{restricted}
\externaldocument[spaces-resolve-]{spaces-resolve}
\externaldocument[formal-defos-]{formal-defos}
\externaldocument[defos-]{defos}
\externaldocument[cotangent-]{cotangent}
\externaldocument[examples-defos-]{examples-defos}
\externaldocument[algebraic-]{algebraic}
\externaldocument[examples-stacks-]{examples-stacks}
\externaldocument[stacks-sheaves-]{stacks-sheaves}
\externaldocument[criteria-]{criteria}
\externaldocument[artin-]{artin}
\externaldocument[quot-]{quot}
\externaldocument[stacks-properties-]{stacks-properties}
\externaldocument[stacks-morphisms-]{stacks-morphisms}
\externaldocument[stacks-limits-]{stacks-limits}
\externaldocument[stacks-cohomology-]{stacks-cohomology}
\externaldocument[stacks-perfect-]{stacks-perfect}
\externaldocument[stacks-introduction-]{stacks-introduction}
\externaldocument[stacks-more-morphisms-]{stacks-more-morphisms}
\externaldocument[stacks-geometry-]{stacks-geometry}
\externaldocument[moduli-]{moduli}
\externaldocument[moduli-curves-]{moduli-curves}
\externaldocument[examples-]{examples}
\externaldocument[exercises-]{exercises}
\externaldocument[guide-]{guide}
\externaldocument[desirables-]{desirables}
\externaldocument[coding-]{coding}
\externaldocument[obsolete-]{obsolete}
\externaldocument[fdl-]{fdl}
\externaldocument[index-]{index}

% Theorem environments.
%
\theoremstyle{plain}
\newtheorem{theorem}[subsection]{Theorem}
\newtheorem{proposition}[subsection]{Proposition}
\newtheorem{lemma}[subsection]{Lemma}

\theoremstyle{definition}
\newtheorem{definition}[subsection]{Definition}
\newtheorem{example}[subsection]{Example}
\newtheorem{exercise}[subsection]{Exercise}
\newtheorem{situation}[subsection]{Situation}

\theoremstyle{remark}
\newtheorem{remark}[subsection]{Remark}
\newtheorem{remarks}[subsection]{Remarks}

\numberwithin{equation}{subsection}

% Macros
%
\def\lim{\mathop{\mathrm{lim}}\nolimits}
\def\colim{\mathop{\mathrm{colim}}\nolimits}
\def\Spec{\mathop{\mathrm{Spec}}}
\def\Hom{\mathop{\mathrm{Hom}}\nolimits}
\def\Ext{\mathop{\mathrm{Ext}}\nolimits}
\def\SheafHom{\mathop{\mathcal{H}\!\mathit{om}}\nolimits}
\def\SheafExt{\mathop{\mathcal{E}\!\mathit{xt}}\nolimits}
\def\Sch{\mathit{Sch}}
\def\Mor{\mathop{\mathrm{Mor}}\nolimits}
\def\Ob{\mathop{\mathrm{Ob}}\nolimits}
\def\Sh{\mathop{\mathit{Sh}}\nolimits}
\def\NL{\mathop{N\!L}\nolimits}
\def\CH{\mathop{\mathrm{CH}}\nolimits}
\def\proetale{{pro\text{-}\acute{e}tale}}
\def\etale{{\acute{e}tale}}
\def\QCoh{\mathit{QCoh}}
\def\Ker{\mathop{\mathrm{Ker}}}
\def\Im{\mathop{\mathrm{Im}}}
\def\Coker{\mathop{\mathrm{Coker}}}
\def\Coim{\mathop{\mathrm{Coim}}}

% Boxtimes
%
\DeclareMathSymbol{\boxtimes}{\mathbin}{AMSa}{"02}

%
% Macros for moduli stacks/spaces
%
\def\QCohstack{\mathcal{QC}\!\mathit{oh}}
\def\Cohstack{\mathcal{C}\!\mathit{oh}}
\def\Spacesstack{\mathcal{S}\!\mathit{paces}}
\def\Quotfunctor{\mathrm{Quot}}
\def\Hilbfunctor{\mathrm{Hilb}}
\def\Curvesstack{\mathcal{C}\!\mathit{urves}}
\def\Polarizedstack{\mathcal{P}\!\mathit{olarized}}
\def\Complexesstack{\mathcal{C}\!\mathit{omplexes}}
% \Pic is the operator that assigns to X its picard group, usage \Pic(X)
% \Picardstack_{X/B} denotes the Picard stack of X over B
% \Picardfunctor_{X/B} denotes the Picard functor of X over B
\def\Pic{\mathop{\mathrm{Pic}}\nolimits}
\def\Picardstack{\mathcal{P}\!\mathit{ic}}
\def\Picardfunctor{\mathrm{Pic}}
\def\Deformationcategory{\mathcal{D}\!\mathit{ef}}
\usepackage{cleveref}
\crefname{originallemma}{original lemma}{original lemmas}
\Crefname{originallemma}{Original lemma}{Original lemmas}
\makeatletter
\newlabel{derived-lemma-leray-acyclicity}{{015E}{source}{}{sourcepointer-1}{}}
\newlabel{derived-lemma-leray-acyclicity@cref}{{[originallemma][0][]015E}{[1][0][]source}{}{}{}}
\makeatother
\crefname{originallemma}{original lemma}{original lemmas}
\Crefname{originallemma}{Original lemma}{Original lemmas}
\makeatletter
\newlabel{hypercovering-lemma-acyclic-hypercover-sheaves}{{01GE}{source}{}{sourcepointer-2}{}}
\newlabel{hypercovering-lemma-acyclic-hypercover-sheaves@cref}{{[originallemma][0][]01GE}{[1][0][]source}{}{}{}}
\makeatother
\crefname{originalreference}{original reference}{original references}
\Crefname{originalreference}{Original reference}{Original references}
\makeatletter
\newlabel{lemma-augmentation-spectral-sequence}{{0D7A}{source}{}{sourcepointer-3}{}}
\newlabel{lemma-augmentation-spectral-sequence@cref}{{[originalreference][0][]0D7A}{[1][0][]source}{}{}{}}
\makeatother
\crefname{originalreference}{original reference}{original references}
\Crefname{originalreference}{Original reference}{Original references}
\makeatletter
\newlabel{lemma-hypercovering-cech-complex}{{0D8F}{source}{}{sourcepointer-4}{}}
\newlabel{lemma-hypercovering-cech-complex@cref}{{[originalreference][0][]0D8F}{[1][0][]source}{}{}{}}
\makeatother
\crefname{originalreference}{original reference}{original references}
\Crefname{originalreference}{Original reference}{Original references}
\makeatletter
\newlabel{lemma-hypercovering-descent-sheaves}{{0D8E}{source}{}{sourcepointer-5}{}}
\newlabel{lemma-hypercovering-descent-sheaves@cref}{{[originalreference][0][]0D8E}{[1][0][]source}{}{}{}}
\makeatother
\crefname{originalreference}{original reference}{original references}
\Crefname{originalreference}{Original reference}{Original references}
\makeatletter
\newlabel{lemma-localize-injective}{{0D88}{source}{}{sourcepointer-6}{}}
\newlabel{lemma-localize-injective@cref}{{[originalreference][0][]0D88}{[1][0][]source}{}{}{}}
\makeatother
\crefname{originalreference}{original reference}{original references}
\Crefname{originalreference}{Original reference}{Original references}
\makeatletter
\newlabel{lemma-restriction-injective-to-component-site}{{09WH}{source}{}{sourcepointer-7}{}}
\newlabel{lemma-restriction-injective-to-component-site@cref}{{[originalreference][0][]09WH}{[1][0][]source}{}{}{}}
\makeatother
\crefname{originalreference}{original reference}{original references}
\Crefname{originalreference}{Original reference}{Original references}
\makeatletter
\newlabel{lemma-simplicial-resolution-augmentation}{{0D78}{source}{}{sourcepointer-8}{}}
\newlabel{lemma-simplicial-resolution-augmentation@cref}{{[originalreference][0][]0D78}{[1][0][]source}{}{}{}}
\makeatother
\crefname{originallemma}{original lemma}{original lemmas}
\Crefname{originallemma}{Original lemma}{Original lemmas}
\makeatletter
\newlabel{more-algebra-lemma-prod-qis-gives-qis}{{0EYX}{source}{}{sourcepointer-9}{}}
\newlabel{more-algebra-lemma-prod-qis-gives-qis@cref}{{[originallemma][0][]0EYX}{[1][0][]source}{}{}{}}
\makeatother
\crefname{originallemma}{original lemma}{original lemmas}
\Crefname{originallemma}{Original lemma}{Original lemmas}
\makeatletter
\newlabel{sites-cohomology-lemma-RHom-into-K-injective}{{0A94}{source}{}{sourcepointer-10}{}}
\newlabel{sites-cohomology-lemma-RHom-into-K-injective@cref}{{[originallemma][0][]0A94}{[1][0][]source}{}{}{}}
\makeatother
\crefname{originallemma}{original lemma}{original lemmas}
\Crefname{originallemma}{Original lemma}{Original lemmas}
\makeatletter
\newlabel{sites-cohomology-lemma-section-RHom-over-U}{{08JA}{source}{}{sourcepointer-11}{}}
\newlabel{sites-cohomology-lemma-section-RHom-over-U@cref}{{[originallemma][0][]08JA}{[1][0][]source}{}{}{}}
\makeatother
\crefname{originalsection}{original section}{original sections}
\Crefname{originalsection}{Original section}{Original sections}
\makeatletter
\newlabel{sites-cohomology-section-internal-hom}{{08J7}{source}{}{sourcepointer-12}{}}
\newlabel{sites-cohomology-section-internal-hom@cref}{{[originalsection][0][]08J7}{[1][0][]source}{}{}{}}
\makeatother
\crefname{originalremark}{original remark}{original remarks}
\Crefname{originalremark}{Original remark}{Original remarks}
\makeatletter
\newlabel{sites-modules-remark-j-shriek-tensor}{{0EYY}{source}{}{sourcepointer-13}{}}
\newlabel{sites-modules-remark-j-shriek-tensor@cref}{{[originalremark][0][]0EYY}{[1][0][]source}{}{}{}}
\makeatother
\title{Derived hypercovering descent: E -> Ra\_*a\textasciicircum{}-1E is an isomorphism in the derived category}
\author{The Stacks Project Authors}
\date{stacks-project:spaces-simplicial.tex:local-snapshot-20261012; GNU FDL 1.2 or later}
\begin{document}
\maketitle
\noindent\textit{Editorial note.} All delivered mathematical text is the author's own, copied from the pinned Stacks Project chapter file sources/stacks-spaces-simplicial/pinned-original-source.tex; no statement or proof step was rewritten, reordered or removed.\par\medskip
\noindent\textit{Editorial note.} The counted proof is the author's complete argument: the author's 'urge the reader to skip' sentence is his own framing of the convergence technicality, not an omission - the double-complex argument that follows is complete and retained verbatim; the degree-zero isomorphism for injective sheaves cites the sheaf-level descent lemma delivered as card L46146.\par\medskip
\noindent\textit{Editorial note.} The label lemma-hypercovering-descent-bounded-abelian names the counted statement, and the proof environment immediately following it in the source is delivered verbatim as its proof.\par\medskip
\noindent\textit{Editorial note.} The counted statement is carried by the author's own lemma environment in the chapter source.\par\medskip
\noindent\textit{Allowed standard prerequisites (editorial).}\begin{itemize}
\item Augmentation spectral sequence for the total site. Same chapter, \textbackslash{}label\{lemma-augmentation-spectral-sequence\}. Line 4393: the spectral sequence whose degeneration on injectives opens the proof.
\item a\_p\textasciicircum{}-1 preserves injectives. Same chapter, \textbackslash{}label\{lemma-localize-injective\}. Line 4396: (a\textasciicircum{}-1I)\_p = a\_p\textasciicircum{}-1I is injective, degenerating the spectral sequence; reused at line 4539 for the K-injective complex.
\item Ra\_*a\textasciicircum{}-1 computed by the hypercovering Cech complex. Same chapter, \textbackslash{}label\{lemma-hypercovering-cech-complex\}. Line 4403: identifies the degenerated complex with SheafHom(s(Z\_\{F(K)\}\textasciicircum{}\#), I).
\item s(Z\_\{F(K)\}\textasciicircum{}\#) is exact in positive degrees for a hypercovering. Hypercoverings chapter, \textbackslash{}label\{hypercovering-lemma-acyclic-hypercover-sheaves\}. Lines 4407-4410: acyclicity of the simplicial presheaf underlying K; applied again at line 4488 to the quasi-isomorphism s(Z\textasciicircum{}\#) -> Z.
\item Sheaf-level hypercovering descent. Same chapter, \textbackslash{}label\{lemma-hypercovering-descent-sheaves\} (own card, statement line 4166). Line 4424: I = a\_*a\textasciicircum{}-1I gives the degree-zero isomorphism for injective sheaves.
\item Leray's acyclicity lemma. Derived Categories chapter, \textbackslash{}label\{derived-lemma-leray-acyclicity\}. Line 4431: derives the bounded-below case from the acyclicity of injectives.
\item Simplicial resolution of the sheaf Z\_U\textasciicircum{}\# under a\textasciicircum{}-1. Same chapter, \textbackslash{}label\{lemma-simplicial-resolution-augmentation\}. Lines 4460-4466: the quasi-isomorphism from the g\_\{p!\}-complex to a\textasciicircum{}-1Z\_U\textasciicircum{}\# opening the unbounded case.
\item RSheafHom against a K-injective complex computes as a product total complex. Cohomology on Sites chapter, Section \textbackslash{}ref\{sites-cohomology-section-internal-hom\}, \textbackslash{}label\{sites-cohomology-lemma-RHom-into-K-injective\}, \textbackslash{}label\{sites-cohomology-lemma-section-RHom-over-U\}. Lines 4472-4483 and 4505-4513: both RGamma sides become Tot\_pi of Hom-term double complexes.
\item j\_! tensor identity for localized sheaves. Modules on Sites chapter, Remark \textbackslash{}ref\{sites-modules-remark-j-shriek-tensor\}. Line 4527: Z\textasciicircum{}\#\_\{K\_p\} tensor Z\_U\textasciicircum{}\# = a\_\{p!\}a\_p\textasciicircum{}-1Z\_U\textasciicircum{}\#, matching the terms of the two double complexes.
\item g\_p\textasciicircum{}-1 preserves K-injectivity. Same chapter, \textbackslash{}label\{lemma-restriction-injective-to-component-site\}. Line 4541: g\_p\textasciicircum{}-1J\textasciicircum{}bullet is K-injective, so the row complexes both compute RHom(a\_p\textasciicircum{}-1Z\_U\textasciicircum{}\#, a\_p\textasciicircum{}-1E).
\item Product of termwise quasi-isomorphisms of double complexes. More on Algebra chapter, \textbackslash{}label\{more-algebra-lemma-prod-qis-gives-qis\}. Line 4550: promotes the termwise quasi-isomorphism A\textasciicircum{}\{p, bullet\} -> B\textasciicircum{}\{p, bullet\} to the product-total quasi-isomorphism.
\end{itemize}\medskip
\ifdefined\hypertarget\hypertarget{sourcepointer-1}{}\fi
\noindent\textit{Source reference.} Original lemma 015E: sources/stacks-chow/tags.txt, source lines 1505--1505. Prerequisite pointer invoked at proof line(s) 4429; it names a standard fact of a sibling chapter that the counted argument uses, and every step of the counted proof built on it is retained verbatim in the delivered body.\par\medskip
\ifdefined\hypertarget\hypertarget{sourcepointer-2}{}\fi
\noindent\textit{Source reference.} Original lemma 01GE: sources/stacks-chow/tags.txt, source lines 1903--1903. Prerequisite pointer invoked at proof line(s) 4414; it names a standard fact of a sibling chapter that the counted argument uses, and every step of the counted proof built on it is retained verbatim in the delivered body.\par\medskip
\ifdefined\hypertarget\hypertarget{sourcepointer-3}{}\fi
\noindent\textit{Source reference.} Original reference 0D7A: sources/stacks-spaces-simplicial/pinned-original-source.tex, source lines 2020--2020. Prerequisite pointer invoked at proof line(s) 4398; it names a standard fact of the same chapter that the counted argument uses, and every step of the counted proof built on it is retained verbatim in the delivered body.\par\medskip
\ifdefined\hypertarget\hypertarget{sourcepointer-4}{}\fi
\noindent\textit{Source reference.} Original reference 0D8F: sources/stacks-spaces-simplicial/pinned-original-source.tex, source lines 4354--4354. Prerequisite pointer invoked at proof line(s) 4409; it names a standard fact of the same chapter that the counted argument uses, and every step of the counted proof built on it is retained verbatim in the delivered body.\par\medskip
\ifdefined\hypertarget\hypertarget{sourcepointer-5}{}\fi
\noindent\textit{Source reference.} Original reference 0D8E: sources/stacks-spaces-simplicial/pinned-original-source.tex, source lines 4167--4167. Prerequisite pointer invoked at proof line(s) 4422; it names a standard fact of the same chapter that the counted argument uses, and every step of the counted proof built on it is retained verbatim in the delivered body.\par\medskip
\ifdefined\hypertarget\hypertarget{sourcepointer-6}{}\fi
\noindent\textit{Source reference.} Original reference 0D88: sources/stacks-spaces-simplicial/pinned-original-source.tex, source lines 3775--3775. Prerequisite pointer invoked at proof line(s) 4401; it names a standard fact of the same chapter that the counted argument uses, and every step of the counted proof built on it is retained verbatim in the delivered body.\par\medskip
\ifdefined\hypertarget\hypertarget{sourcepointer-7}{}\fi
\noindent\textit{Source reference.} Original reference 09WH: sources/stacks-spaces-simplicial/pinned-original-source.tex, source lines 1086--1086. Prerequisite pointer invoked at proof line(s) 4539; it names a standard fact of the same chapter that the counted argument uses, and every step of the counted proof built on it is retained verbatim in the delivered body.\par\medskip
\ifdefined\hypertarget\hypertarget{sourcepointer-8}{}\fi
\noindent\textit{Source reference.} Original reference 0D78: sources/stacks-spaces-simplicial/pinned-original-source.tex, source lines 1929--1929. Prerequisite pointer invoked at proof line(s) 4459; it names a standard fact of the same chapter that the counted argument uses, and every step of the counted proof built on it is retained verbatim in the delivered body.\par\medskip
\ifdefined\hypertarget\hypertarget{sourcepointer-9}{}\fi
\noindent\textit{Source reference.} Original lemma 0EYX: sources/stacks-chow/tags.txt, source lines 18707--18707. Prerequisite pointer invoked at proof line(s) 4553; it names a standard fact of a sibling chapter that the counted argument uses, and every step of the counted proof built on it is retained verbatim in the delivered body.\par\medskip
\ifdefined\hypertarget\hypertarget{sourcepointer-10}{}\fi
\noindent\textit{Source reference.} Original lemma 0A94: sources/stacks-chow/tags.txt, source lines 12937--12937. Prerequisite pointer invoked at proof line(s) 4490; it names a standard fact of a sibling chapter that the counted argument uses, and every step of the counted proof built on it is retained verbatim in the delivered body.\par\medskip
\ifdefined\hypertarget\hypertarget{sourcepointer-11}{}\fi
\noindent\textit{Source reference.} Original lemma 08JA: sources/stacks-chow/tags.txt, source lines 10843--10843. Prerequisite pointer invoked at proof line(s) 4491; it names a standard fact of a sibling chapter that the counted argument uses, and every step of the counted proof built on it is retained verbatim in the delivered body.\par\medskip
\ifdefined\hypertarget\hypertarget{sourcepointer-12}{}\fi
\noindent\textit{Source reference.} Original section 08J7: sources/stacks-chow/tags.txt, source lines 10840--10840. Prerequisite pointer invoked at proof line(s) 4480; it names a standard fact of a sibling chapter that the counted argument uses, and every step of the counted proof built on it is retained verbatim in the delivered body.\par\medskip
\ifdefined\hypertarget\hypertarget{sourcepointer-13}{}\fi
\noindent\textit{Source reference.} Original remark 0EYY: sources/stacks-chow/tags.txt, source lines 18708--18708. Prerequisite pointer invoked at proof line(s) 4528; it names a standard fact of a sibling chapter that the counted argument uses, and every step of the counted proof built on it is retained verbatim in the delivered body.\par\medskip
\begin{lemma}
\label{lemma-hypercovering-descent-bounded-abelian}
Let $\mathcal{C}$ be a site with equalizers and fibre products.
Let $K$ be a hypercovering. For
$E \in D(\mathcal{C})$ the map
$$
E \longrightarrow Ra_*a^{-1}E
$$
is an isomorphism.
\end{lemma}

\begin{proof}
First, let $\mathcal{I}$ be an injective abelian sheaf on $\mathcal{C}$.
Then the spectral sequence of
Lemma \ref{lemma-augmentation-spectral-sequence}
for the sheaf $a^{-1}\mathcal{I}$ degenerates as
$(a^{-1}\mathcal{I})_p = a_p^{-1}\mathcal{I}$
is injective by Lemma \ref{lemma-localize-injective}.
Thus the complex
$$
a_{0, *}a_0^{-1}\mathcal{I} \to
a_{1, *}a_1^{-1}\mathcal{I} \to
a_{2, *}a_2^{-1}\mathcal{I} \to\ldots
$$
computes $Ra_*a^{-1}\mathcal{I}$. By
Lemma \ref{lemma-hypercovering-cech-complex}
this is equal to the complex
$\SheafHom(s(\mathbf{Z}_{F(K)}^\#), \mathcal{I})$.
Because $K$ is a hypercovering, we see that
$s(\mathbf{Z}_{F(K)}^\#)$ is exact in degrees $> 0$ by
Hypercoverings, Lemma \ref{hypercovering-lemma-acyclic-hypercover-sheaves}
applied to the simplicial presheaf $F(K)$.
Since $\mathcal{I}$ is injective, the functor $\SheafHom(-, \mathcal{I})$
is exact and we conclude that
$\SheafHom(s(\mathbf{Z}_{F(K)}^\#), \mathcal{I})$
is exact in positive degrees. We conclude that
$R^pa_*a^{-1}\mathcal{I} = 0$ for $p > 0$.
On the other hand, we have $\mathcal{I} = a_*a^{-1}\mathcal{I}$
by Lemma \ref{lemma-hypercovering-descent-sheaves}.

\medskip\noindent
Bounded case. Let $E \in D^+(\mathcal{C})$.
Choose a bounded below complex $\mathcal{I}^\bullet$ of injectives
representing $E$. By the result of the first paragraph and
Leray's acyclicity lemma
(Derived Categories, Lemma \ref{derived-lemma-leray-acyclicity})
$Ra_*a^{-1}\mathcal{I}^\bullet$
is computed by the complex
$a_*a^{-1}\mathcal{I}^\bullet = \mathcal{I}^\bullet$
and we conclude the lemma is true in this case.

\medskip\noindent
Unbounded case. We urge the reader to skip this, since the argument
is the same as above, except that we use explicit representation
by double complexes to get around convergence issues.
Let $E \in D(\mathcal{C})$.
To show the map $E \to Ra_*a^{-1}E$ is an isomorphism,
it suffices to show for every object $U$ of $\mathcal{C}$ that
$$
R\Gamma(U, E) = R\Gamma(U, Ra_*a^{-1}E)
$$
We will compute both sides and show the map $E \to Ra_*a^{-1}E$
induces an isomorphism. Choose a K-injective
complex $\mathcal{I}^\bullet$ representing $E$. Choose a quasi-isomorphism
$a^{-1}\mathcal{I}^\bullet \to \mathcal{J}^\bullet$
for some K-injective complex $\mathcal{J}^\bullet$ on
$(\mathcal{C}/K)_{total}$. We have
$$
R\Gamma(U, E) = R\Hom(\mathbf{Z}_U^\#, E)
$$
and
$$
R\Gamma(U, Ra_*a^{-1}E) = R\Hom(\mathbf{Z}_U^\#, Ra_*a^{-1}E) =
R\Hom(a^{-1}\mathbf{Z}_U^\#, a^{-1}E)
$$
By Lemma \ref{lemma-simplicial-resolution-augmentation}
we have a quasi-isomorphism
$$
\Big(\ldots \to
g_{2!}(a_2^{-1}\mathbf{Z}_U^\#) \to
g_{1!}(a_1^{-1}\mathbf{Z}_U^\#) \to
g_{0!}(a_0^{-1}\mathbf{Z}_U^\#)\Big)
\longrightarrow
a^{-1}\mathbf{Z}_U^\#
$$
Hence $R\Hom(a^{-1}\mathbf{Z}_U^\#, a^{-1}E)$ is equal to
$$
R\Gamma((\mathcal{C}/K)_{total},
R\SheafHom(
\ldots \to
g_{2!}(a_2^{-1}\mathbf{Z}_U^\#) \to
g_{1!}(a_1^{-1}\mathbf{Z}_U^\#) \to
g_{0!}(a_0^{-1}\mathbf{Z}_U^\#),
\mathcal{J}^\bullet))
$$
By the construction in Cohomology on Sites, Section
\ref{sites-cohomology-section-internal-hom}
and since $\mathcal{J}^\bullet$ is K-injective, we see that
this is represented by the complex of abelian groups with terms
$$
\prod\nolimits_{p + q = n}
\Hom(g_{p!}(a_p^{-1}\mathbf{Z}_U^\#), \mathcal{J}^q) =
\prod\nolimits_{p + q = n}
\Hom(a_p^{-1}\mathbf{Z}_U^\#, g_p^{-1}\mathcal{J}^q)
$$
See Cohomology on Sites, Lemmas
\ref{sites-cohomology-lemma-RHom-into-K-injective} and
\ref{sites-cohomology-lemma-section-RHom-over-U} for more information.
Thus we find that $R\Gamma(U, Ra_*a^{-1}E)$ is computed by
the product total complex $\text{Tot}_\pi(B^{\bullet, \bullet})$
with $B^{p, q} = \Hom(a_p^{-1}\mathbf{Z}_U^\#, g_p^{-1}\mathcal{J}^q)$.
For the other side we argue similarly. First we note that
$$
s(\mathbf{Z}_{F(K)}^\#) \longrightarrow \mathbf{Z}
$$
is a quasi-isomorphism of complexes on $\mathcal{C}$
by Hypercoverings, Lemma \ref{hypercovering-lemma-acyclic-hypercover-sheaves}.
Since $\mathbf{Z}_U^\#$ is a flat sheaf of $\mathbf{Z}$-modules
we see that
$$
s(\mathbf{Z}_{F(K)}^\#) \otimes_\mathbf{Z} \mathbf{Z}_U^\#
\longrightarrow
\mathbf{Z}_U^\#
$$
is a quasi-isomorphism. Therefore
$R\Hom(\mathbf{Z}_U^\#, E)$ is equal to
$$
R\Gamma(\mathcal{C}, R\SheafHom(
s(\mathbf{Z}_{F(K)}^\#) \otimes_\mathbf{Z} \mathbf{Z}_U^\#,
\mathcal{I}^\bullet))
$$
By the construction of $R\SheafHom$ and since $\mathcal{I}^\bullet$
is K-injective, this is represented by the complex of abelian groups
with terms
$$
\prod\nolimits_{p + q = n}
\Hom(\mathbf{Z}^\#_{K_p} \otimes_\mathbf{Z} \mathbf{Z}_U^\#, \mathcal{I}^q)
=
\prod\nolimits_{p + q = n}
\Hom(a_p^{-1}\mathbf{Z}_U^\#, a_p^{-1}\mathcal{I}^q)
$$
The equality of terms follows from the fact that
$\mathbf{Z}^\#_{K_p} \otimes_\mathbf{Z} \mathbf{Z}_U^\# =
a_{p!}a_p^{-1}\mathbf{Z}_U^\#$ by Modules on Sites,
Remark \ref{sites-modules-remark-j-shriek-tensor}.
Thus we find that $R\Gamma(U, E)$ is computed by
the product total complex $\text{Tot}_\pi(A^{\bullet, \bullet})$
with $A^{p, q} = \Hom(a_p^{-1}\mathbf{Z}_U^\#, a_p^{-1}\mathcal{I}^q)$.

\medskip\noindent
Since $\mathcal{I}^\bullet$ is K-injective we see that
$a_p^{-1}\mathcal{I}^\bullet$ is K-injective, see
Lemma \ref{lemma-localize-injective}.
Since $\mathcal{J}^\bullet$ is K-injective we see that
$g_p^{-1}\mathcal{J}^\bullet$ is K-injective, see
Lemma \ref{lemma-restriction-injective-to-component-site}.
Both represent the object $a_p^{-1}E$.
Hence for every $p \geq 0$ the map of complexes
$$
A^{p, \bullet} = \Hom(a_p^{-1}\mathbf{Z}_U^\#, a_p^{-1}\mathcal{I}^\bullet)
\longrightarrow
\Hom(a_p^{-1}\mathbf{Z}_U^\#, g_p^{-1}\mathcal{J}^\bullet) = B^{p, \bullet}
$$
induced by $g_p^{-1}$ applied to the given map
$a^{-1}\mathcal{I}^\bullet \to \mathcal{J}^\bullet$
is a quasi-isomorphisms as these complexes both compute
$$
R\Hom(a_p^{-1}\mathbf{Z}_U^\#, a_p^{-1}E)
$$
By More on Algebra, Lemma \ref{more-algebra-lemma-prod-qis-gives-qis}
we conclude that the right vertical arrow in the commutative diagram
$$
\xymatrix{
R\Gamma(U, E) \ar[r] \ar[d] &
\text{Tot}_\pi(A^{\bullet, \bullet}) \ar[d] \\
R\Gamma(U, Ra_*a^{-1}E) \ar[r] &
\text{Tot}_\pi(B^{\bullet, \bullet})
}
$$
is a quasi-isomorphism. Since we saw above that the horizontal arrows
are quasi-isomorphisms, so is the left vertical arrow.
\end{proof}

\end{document}
